\section{Vektorer og Matricer}
\begin{minipage}{0.35\textwidth}
    \[
        \mathbb{F}^{3 \times 1} =
        \left\{
            \begin{bmatrix}
                a_{1} \\
                a_{2} \\
                a_{3}
                \end{bmatrix}
            \mid a_{1}, a_{2}, a_{3} \in \mathbb{F}
        \right\}
    \]
\end{minipage}
\hfill
\begin{minipage}{0.55\textwidth}
    \[
        \mathbb{F}^{1 \times 3} =
        \left\{
            \begin{bmatrix}
                a_{1} & a_{2} & a_{3}
            \end{bmatrix}
            \mid a_{1}, a_{2}, a_{3} \in \mathbb{F}
        \right\}
    \]
\end{minipage}


\subsection{Addition og Subtraktion med matrix}
Kun defineret hvis begge matricer har samme størrelse.

\[
    \begin{bmatrix}
        a_{11} & a_{12} \\
        a_{21} & a_{22}
    \end{bmatrix}
    +
    \begin{bmatrix}
        b_{11} & b_{12} \\
        b_{21} & b_{22}
    \end{bmatrix}
    =
    \begin{bmatrix}
        a_{11} + b_{11} & a_{12} + b_{12} \\
        a_{21} + b_{21} & a_{22} + b_{22}
    \end{bmatrix}
\]

\[
    \begin{bmatrix}
        a_{11} & a_{12} \\
        a_{21} & a_{22}
    \end{bmatrix}
    -
    \begin{bmatrix}
        b_{11} & b_{12} \\
        b_{21} & b_{22}
    \end{bmatrix}
    =
    \begin{bmatrix}
        a_{11} - b_{11} & a_{12} - b_{12} \\
        a_{21} - b_{21} & a_{22} - b_{22}
    \end{bmatrix}
\]

\subsection{Multiplikation af vektor og skalar, matrix og skalar}
\begin{minipage}{0.25\textwidth}
    \[
        k \cdot
        \begin{bmatrix}
            a_{1} \\
            a_{2} \\
            a_{3}
        \end{bmatrix}
        =
        \begin{bmatrix}
            k \cdot a_{1} \\
            k \cdot a_{2} \\
            k \cdot a_{3}
        \end{bmatrix}
    \]
\end{minipage}
\hfill
\begin{minipage}{0.55\textwidth}
    \[
        k \cdot
        \begin{bmatrix}
            a_{11} & a_{12} & a_{13} \\
            a_{21} & a_{22} & a_{23} \\
            a_{31} & a_{32} & a_{33}
        \end{bmatrix}
        =
        \begin{bmatrix}
            k \cdot a_{11} & k \cdot a_{12} & k \cdot a_{13} \\
            k \cdot a_{21} & k \cdot a_{22} & k \cdot a_{23} \\
            k \cdot a_{31} & k \cdot a_{32} & k \cdot a_{33}
        \end{bmatrix}
    \]
\end{minipage}


\subsection{Multiplikation af matrix og vektor}
Kun defineret hvis antallet af kolonner i matricen er lig med antallet af rækker i vektoren. Resultatet bliver en vektor med samme antal rækker som matricen.

\[
    \begin{bmatrix}
        a_{11} & a_{12} & a_{13} \\
        a_{21} & a_{22} & a_{23}
    \end{bmatrix}
    \cdot
    \begin{bmatrix}
        b_{1} \\
        b_{2} \\
        b_{3}
    \end{bmatrix}
    =
    \begin{bmatrix}
        a_{11}b_{1} + a_{12}b_{2} + a_{13}b_{3} \\
        a_{21}b_{1} + a_{22}b_{2} + a_{23}b_{3}
    \end{bmatrix}
\]

\subsection{Multiplikation af matrix og matrix}
Kun defineret hvis antallet af søjler i den første matrix er lig med antallet af rækker i den anden matrix. Resultatet bliver en matrix med antal rækker fra den første matrix og antal kolonner fra den anden matrix.

\[
    \begin{bmatrix}
        a_{11} & a_{12} \\
        a_{21} & a_{22}
    \end{bmatrix}
    \cdot
    \begin{bmatrix}
        b_{11} & b_{12} & b_{13} \\
        b_{21} & b_{22} & b_{23}
    \end{bmatrix}
    =
    \begin{bmatrix}
        (a_{11}b_{11} + a_{12}b_{21}) & (a_{11}b_{12} + a_{12}b_{22}) & (a_{11}b_{13} + a_{12}b_{23}) \\
        (a_{21}b_{11} + a_{22}b_{21}) & (a_{21}b_{12} + a_{22}b_{22}) & (a_{21}b_{13} + a_{22}b_{23})
    \end{bmatrix}
\]


\bemærk{$A \cdot B \neq B \cdot A$}

\subsection{Algebra}
\subsubsection{Vektorer}
\begin{align*}
    (\mathbf{u} + \mathbf{v}) + \mathbf{w} &= \mathbf{u} + (\mathbf{v} + \mathbf{w}) \\
    \mathbf{u} + \mathbf{v} &= \mathbf{v} + \mathbf{u} \\
    c \cdot (d \cdot \mathbf{u}) &= (c \cdot d) \cdot \mathbf{u} \\
    c \cdot (\mathbf{u} + \mathbf{v}) &= c \cdot \mathbf{u} + c \cdot \mathbf{v} \\
    (c + d) \cdot \mathbf{u} &= c \cdot \mathbf{u} + d \cdot \mathbf{u}
\end{align*}
\kilde{Sætning 7.1.1}

\subsubsection{Matricer}
For $\mathbf{A} \in \mathbb{F}^{m \times n}$, $\mathbf{B} \in \mathbb{F}^{n \times l}$ og $\mathbf{C} \in \mathbb{F}^{l \times k}$ gælder:


\begin{align*}
    \mathbf{A}_1 + \mathbf{A}_2 &= \mathbf{A}_2 + \mathbf{A}_1 \\
    (\mathbf{A}_1 + \mathbf{A}_2) + \mathbf{A}_3 &= \mathbf{A}_1 + (\mathbf{A}_2 + \mathbf{A}_3) \\
    \mathbf{A} \cdot (\mathbf{B} \cdot \mathbf{C}) &= (\mathbf{A} \cdot \mathbf{B}) \cdot \mathbf{C} \\
    \mathbf{A} \cdot (\mathbf{B}_1 + \mathbf{B}_2) &= \mathbf{A} \cdot \mathbf{B}_1 + \mathbf{A} \cdot \mathbf{B}_2 \\
    (\mathbf{A}_1 + \mathbf{A}_2) \cdot \mathbf{B} &= \mathbf{A}_1 \cdot \mathbf{B} + \mathbf{A}_2 \cdot \mathbf{B} 
\end{align*}
\kilde{Sætning 7.2.1}

\subsection{Transponering af vektorer}
\[
    \begin{bmatrix}
        a_{1} & a_{2} & a_{3}
    \end{bmatrix}^{T}
    =
    \begin{bmatrix}
        a_{1} \\
        a_{2} \\
        a_{3}
    \end{bmatrix}
\]

\subsection{Transponering af matricer}
\[
    \begin{bmatrix}
        a_{11} & a_{12} & a_{13} \\
        a_{21} & a_{22} & a_{23} \\
        a_{31} & a_{32} & a_{33}
    \end{bmatrix}^{T}
    =
    \begin{bmatrix}
        a_{11} & a_{21} & a_{31} \\
        a_{12} & a_{22} & a_{32} \\
        a_{13} & a_{23} & a_{33}
    \end{bmatrix}
\]

For $\mathbf{A} \in \mathbb{F}^{m \times n}$ og $\mathbf{B} \in \mathbb{F}^{n \times l}$ gælder:

\begin{itemize}
    \item $(\mathbf{A}^T)^T=\mathbf{A}$
    \item $(\mathbf{A}_1+\mathbf{A}_2)^T=\mathbf{A}_1^T+\mathbf{A}_2^T$
    \item $(\mathbf{A} \cdot \mathbf{B})^T = \mathbf{B}^T \cdot \mathbf{A}^T $ \hfill \bemærk{Rækkefølgen vendes}
\end{itemize}

\kilde{Sætning 7.2.2}

\subsection{Inverse matricer}
\begin{itemize}
    \item Kun defineret for kvadratiske matricer.
    \item $A$ har en invers $\iff$ $\rho(A) =$ antallet af rækker i $A$.
    \kilde{Lemma 7.3.1}
    \item $A \cdot A^{-1}=I_n$ \hfill \kilde{Korollar 7.3.4}
    \item $A$ har en invers $\iff \det(A) \neq 0$ \hfill \kilde{Korollar 8.2.4}
\end{itemize}



\subsubsection{Invers ved rækkeoperationer}
For en kvadratisk matrix $A$ kan $A^{-1}$ findes ved at bruge rækkeoperationer:
\[
\bigl[\,A \mid I\,\bigr]
\;\xrightarrow{\text{rækkeoperationer}}\;
\bigl[\,I \mid A^{-1}\,\bigr].
\]
Hvis venstre side efter reduktion ikke ender som identitetsmatricen, eksisterer $A^{-1}$ ikke (matricen er singulær).

\[
    \mathbf{A}=
    \begin{bmatrix}
        a_{11} & a_{12} \\
        a_{21} & a_{22}
    \end{bmatrix}
    \implies
    \begin{bmatrix}
        a_{11} & a_{12} & 1 & 0 \\
        a_{21} & a_{22} & 0 & 1
    \end{bmatrix}
    \xrightarrow{\text{række op.}}
    \begin{bmatrix}
        1 & 0 & x & y \\
        0 & 1 & z & w
    \end{bmatrix}
    \implies
    \mathbf{A}^{-1}=
    \begin{bmatrix}
        x & y \\
        z & w
    \end{bmatrix}
\]

\kilde{Sætning 7.3.2}

\subsubsection{Invers ved detminant}
\[
\begin{bmatrix}
    a & b \\
    c & d
\end{bmatrix}
^{-1}
=
\frac{1}{\det(A)} \cdot
\begin{bmatrix}
    d & -b \\
    -c & a
\end{bmatrix}
\kildetag{Sætning 8.1.1}
\]

\subsection{Lineær afhængighed}
Hvis en eller flere vektorer i et sæt kan skrives som en lineær kombination af de andre vektorer i sættet, er vektorerne lineært afhængige. Ellers er de lineært uafhængige.

\subsubsection{Determinant}
Hvis $\det(A) = 0$ er søjlerne i $\mathbf{A}$ lineært afhængige. Ellers uafhængige.

\kilde{Korollar 8.2.5}

\subsubsection{Samme rang som vektorer}
Hvis matricen skabt af vekotrerne som søjler har samme rang som antallet af vektorer, er de lineært uafhængige. Ellers er de lineært afhængige.

\kilde{Sætning 7.1.3}

\newpage
\section{Determinanter}
Kun defineret for kvadratiske matricer. \hfill $\det([a]) = a$

\subsection{Determinant af 2×2 og 3×3 matrix}
\begin{align*}
    \det
    \begin{bmatrix}
        a & b \\
        c & d
    \end{bmatrix}
    &=
    ad - bc \\
    \det
    \begin{bmatrix}
        a & b & c \\
        d & e & f \\
        g & h & i
    \end{bmatrix}
    &=
    aei + bfg + cdh - ceg - bdi - afh
    \kildetag{Sarrus' regel}
\end{align*}

\subsection{Determinant af større matricer (Laplace udvikling)}
\[
    \det
    \begin{bmatrix}
        a_{11} & a_{12} & a_{13} \\
        a_{21} & a_{22} & a_{23} \\
        a_{31} & a_{32} & a_{33}
    \end{bmatrix}
    =
    a_{11} \cdot \det \begin{bmatrix}
        a_{22} & a_{23} \\
        a_{32} & a_{33}
    \end{bmatrix} - 
    a_{12} \cdot \det \begin{bmatrix}
        a_{21} & a_{23} \\
        a_{31} & a_{33}
    \end{bmatrix} +
    a_{13} \cdot \det \begin{bmatrix}
        a_{21} & a_{22} \\
        a_{31} & a_{32}
    \end{bmatrix}
\]


General formel for en valgfri $i$'te række:

\[
    \det(\mathbf{A}) = \sum_{j=1}^{n} (-1)^{i+j} \cdot a_{ij} \cdot \det(\mathbf{A}(i;j))
\]

Generel formel for en valgfri $j$'te kolonne:

\[
    \det(\mathbf{A}) = \sum_{i=1}^{n} (-1)^{i+j} \cdot a_{ij} \cdot \det(\mathbf{A}(i;j))
\]

\kilde{Sætning 8.2.1}

\bemærkBig{
$\mathbf{A}(i;j)$ er \emph{minoren} (matricen) dannet ved at fjerne den $i$'te række og $j$'te kolonne fra matrix $\mathbf{A}$.

Enhver række eller kolonne kan vælges - Brugbart hvis en række eller kolonne indeholder mange nuller.
}

Man tager altså hvert element i den valgte række eller kolonne, ganger det med determinanten af minoren dannet ved at fjerne den række og kolonne, som elementet befinder sig i, og ganger med $(-1)^{i+j}$ for at tage højde for fortegnet.
\renewcommand{\arraystretch}{1.3}
\begin{center}
    \textbf{Fortegns skema for en $n \times n$ matrix}
    
    \begin{tabular}{|c|c|c|c|}
        \hline
        $+a_{11}$ & $-a_{12}$ & $+a_{13}$ & $\cdots$ \\
        \hline
        $-a_{21}$ & $+a_{22}$ & $-a_{23}$ & $\cdots$ \\
        \hline
        $+a_{31}$ & $-a_{32}$ & $+a_{33}$ & $\cdots$ \\
        \hline
        $\vdots$  & $\vdots$  & $\vdots$  & $\ddots$ \\
    \hline
    \end{tabular}
\end{center}



\subsection{Egenskaber ved determinanter}
Hvis en række eller kolonne er nul, er determinanten $0$. 

\kilde{Lemma 8.1.4}

\

Hvis to rækker eller kolonner er ens, er determinanten $0$.

\kilde{Proposition 8.3.5}

\

Hvis to rækker eller kolonner byttes om, skifter determinanten fortegn.

\kilde{Sætning 8.3.6}

\

Hvis en række eller kolonne ganges med en skalar $k$, ganges determinanten med $k$.

\kilde{Korollar 8.3.2}

\

Hvis et skalar af en række eller kolonne lægges til en anden række eller kolonne, ændres determinanten ikke.

\kilde{Sætning 8.3.7}

\

$\det(A) \ne 0 \iff$ det tilsvarende lineære ligningssystem's eneste løsning er nulvektoren. 

\kilde{Korollar 8.2.6}

\begin{align*}
    \det (\mathbf{A}^{-1} ) &= \frac{1}{\det(\mathbf{A})} \kildetag{Sætning 8.1.1} \\
    \text{For diagonalmatricer: } \det(\mathbf{D}) &= \prod_{i=1}^{n} d_{ii} \kildetag{Proposition 8.1.2} \\
    \text{For trekantsmatricer: } \det(\mathbf{T}) &= \prod_{i=1}^{n} t_{ii} \kildetag{Sætning 8.1.(3/5)} \\
    \det(A^T) &= \det(A) \kildetag{Korollar 8.2.2} \\
    \det (\mathbf{A} \cdot \mathbf{B}) &= \det(\mathbf{A}) \cdot \det(\mathbf{B}) \kildetag{Sætning 8.2.3}
\end{align*}